\chapter{並列化の設計と実装}
\label{ch:implementation}

本章では、第2章で述べた Eppstein アルゴリズムを並列化するための設計と実装を述べる。
実装はどれも同じ骨格を共有し、変える点を一つに絞った変種の集まりとして構成した。
バッチの配り方だけを変えた三つの変種、グラフを worker へ届ける方式だけを変えた二つの変種であり、性能差の原因を一つの設計判断に帰着できるようにするためである。
各変種の実測結果は第5章で示す。

\section{最外層の独立性とカウント版への限定}

並列化の足場は、Eppstein アルゴリズムの最外層が持つ独立性にある。
このアルゴリズムは縮退順序に沿って各頂点 $v_i$ を一度ずつ根に取り、その頂点を含む極大クリークだけを担当する部分問題を解く。
頂点 $v_i$ の部分問題は、固定されたグラフと縮退順序だけに依存し、ほかの頂点の部分問題の結果には一切依存しない。
したがって最外層のループは、そのまま互いに独立なタスクの集まりとして切り出せる。

この独立性を、それ以上の工夫を足さずにそのまま並列化の単位とした。
再帰の内側（ピボット選択と候補集合の絞り込み）には手を入れず、外側の反復だけを複数のプロセスに振り分ける。
内側にも並列性はあるが、部分問題どうしがすでに独立している以上、まず最外層だけを配る素朴な構成で固定費と負荷分布を測り、そこで見えた限界に応じて次を決める方針を採った。

並列化の対象は、極大クリークの個数を数える\textbf{カウント版}に限定した。
列挙版は極大クリークの集合そのものを出力するが、soc-sign-epinions では出力が2千万個を超える。
これを複数プロセスから集約するには、各 worker が見つけたクリークをプロセス間で転送してまとめる必要があり、列挙結果の転送と集約という別の重い問題が主役になってしまう。
本研究が測りたいのは列挙計算そのものの並列化なので、出力は各 worker がローカルに数えた整数の総和で済むカウント版を対象とし、転送コストを個数の足し合わせだけに切り離した。

\section{プロセスプールによる基本実装}

基本実装は、Python 標準ライブラリの \verb|multiprocessing.Pool| による素朴なプロセスプールである。
縮退順序の計算は親プロセスで逐次に一度だけ行い、その順序を連続する \verb|BATCH_SIZE| 頂点ずつのバッチに区切って worker へ配る。
各 worker は担当バッチ内の頂点それぞれについて、近傍のうち縮退順序で後ろにあるものを候補集合、前にあるものを除外集合として、ピボット再帰で部分木の極大クリーク数を数える。

worker へのデータ配布には、プールの初期化関数を使う。
\verb|Pool| の \verb|initializer| でグラフと縮退順序を各 worker のモジュール大域変数に一度だけ格納し、以降のタスクはバッチの範囲だけを受け取る。
こうしないと、グラフをタスクごとに pickle 化して送り直すことになり、転送コストがタスク数だけ積み上がる。
macOS のプロセス起動方式は既定で spawn なので、この初期データはプール起動時に pickle 化されてパイプ経由で各 worker へ転送される。
グラフ全体を worker ごとに一度コピーするこの設計は、後で固定費として効いてくる（第5章）が、まずは最も素朴な形を基準線として測るために選んだ。

負荷分散は \verb|imap_unordered| の動的配分に任せた。
バッチは完成した順に結果を回収でき、速く終わった worker から次のバッチを引く。
静的に「どの worker がどのバッチを担当するか」を決め打ちせず、実行時のタスク取得だけで均す狙いである。

\begin{verbatim}
with Pool(workers, initializer=_init_worker,
          initargs=(graph, ordering)) as pool:
    pool.map(_noop, range(workers * 4), chunksize=1)  # 起動バリア
    for pid, elapsed, sub_count, largest in pool.imap_unordered(
            _count_batch_timed, batches):
        count += sub_count
\end{verbatim}

\verb|BATCH_SIZE| は、動的配分が働くだけの細かさと、タスク発行のオーバーヘッドが無視できるだけの粗さの折り合いで決まる。
バッチが大きすぎると、重いバッチを引いた worker が一台だけ長く走り、ほかが待つ。
逆に細かすぎれば、タスクの受け渡し回数が増える。
この粒度が並列化の効きを左右することは第5章で示す。

\section{バッチ配分の三つの戦略}

基本実装のバッチは縮退順序の先頭から連続に区切るが、この区切り方には偏りの問題がある。
負荷分布の測定（第5章）で、重い頂点は縮退順序の末尾に密集することがわかっている。
疎な外周から頂点を剥がしていく縮退順序では、密なコアが必然的に末尾へ回るためである。
連続区切りだと重い頂点が同じ末尾のバッチにまとまり、そのバッチを引いた worker だけが突出して長く走る。

この偏りにバッチの作り方だけで対処するため、区切り方の異なる三つの戦略を用意した。
枠組みはどれも「縮退順序の各頂点を根とする部分問題を、バッチにまとめて worker に配る」で共通し、違うのはバッチの中身の決め方だけである。

\begin{itemize}
\item \textbf{block}（基本実装）：順序の先頭から連続に区切る。重いバッチが末尾にまとまり、最後に配られる。
\item \textbf{reversed}（末尾逆順）：順序を逆にしてから連続に区切る。重いバッチを最初に配り、終盤で重いバッチを引いて長く待つ事態を防ぐ。
\item \textbf{interleave}（ストライド）：バッチ数を $K$ として、バッチ $k$ に順序上の頂点 $k, k+K, k+2K, \dots$ を入れる。どのバッチにも軽い頂点と重い頂点が混ざり、バッチの重さ自体がほぼ均一になる。
\end{itemize}

\begin{verbatim}
# block:      batches = [range(i, min(i + b, n)) for i in range(0, n, b)]
# reversed:   idx = list(range(n))[::-1]
#             batches = [idx[i:i+b] for i in range(0, n, b)]
# interleave: K = -(-n // b)  # バッチ数
#             batches = [range(k, n, K) for k in range(K)]
\end{verbatim}

reversed と interleave は、同じ偏りに別の角度から当たっている。
reversed は重い順に配って終盤の待ちを消すが、最重のバッチ1個より速くは終われない。
interleave は各バッチの重さをならして、最重のバッチ自体を軽くする。
どちらが効くか、そもそも静的な配り方で偏りを救えるのかは第5章で検証する。

\section{グラフを worker へ届ける二つの方式}

基本実装の固定費は、グラフを spawn で各 worker へ pickle 転送するところに集中する。
この固定費を削るために、グラフの届け方だけを変えた二つの変種を用意した。
どちらもバッチの作り方は block と同一にし、block との差がすべて配布方式に帰着するようにしてある。

\subsection*{fork による copy-on-write 継承}

一つ目は、プロセス起動方式を spawn から fork に変える変種である。
\verb|multiprocessing.get_context("fork")| で worker を起動すると、子プロセスは親のメモリを copy-on-write で継承する。
グラフは親のアドレス空間にすでにあるので、pickle 化もパイプ転送も worker 側の再構築も起きず、起動コストは fork 自体まで縮むはずである。

macOS が既定で fork を使わないのは、スレッドを持つプロセスを fork すると、ロックを保持したままのスレッドが子には複製されず、子がデッドロックしうるためである。
本実装は fork を呼ぶ時点で単一スレッドの純 Python であり、この危険がないので fork を安全に使える。
既定を外す判断の根拠がこの単一スレッド性にあることを、変種の docstring にも明記した。

\subsection*{shared\_memory による共有配列}

二つ目は、起動方式は spawn のまま、グラフを共有メモリに置く変種である。
親プロセスがグラフを CSR 形式の int64 配列4本（頂点 ID、隣接リストの開始位置、隣接先、縮退順序）にまとめ、\verb|multiprocessing.shared_memory| のブロックに書き込む。
各 worker は転送を受け取る代わりに、ブロックへ名前で attach する。

ただし、この変種でも worker 側の構造構築は消えない。
再帰カーネルは Python の集合演算（\verb|set|）に依存し、\verb|set| は共有メモリに置けないためである。
共有メモリに置けるのは連続したバイト列だけで、ハッシュテーブルである \verb|set| の内部構造は載らない。
そのため各 worker は、attach した共有配列から自分用の集合の辞書を起動時に組み立て直す。

\begin{verbatim}
graph = {ids[i]: {ids[j] for j in indices[indptr[i]:indptr[i + 1]]}
         for i in range(len(ids))}   # 共有 CSR 配列から set を再構築
\end{verbatim}

この設計により、shm 変種が消すのは pickle 化とパイプ転送だけになる。
spawn の起動コストを「インタプリタ起動」「転送」「構造構築」の三つに分けたとき、転送だけを消すのが shm、三つとも消すのが fork という対比になる。
shm を単独の高速化策としてではなく、spawn の固定費の内訳を切り分ける測定器として位置づけたのはこのためである。

\section{フェーズ分解と worker 稼働時間の計測}

並列化の効き方を診断するために、1回の実行を四つの区間に分けて壁時計時間を測る計測器を実装した。
区間は、グラフの読込または生成（load）、縮退順序の計算（ordering、逐次）、プールの生成と worker 初期化（startup）、並列本体（compute）である。
load と ordering は逐次項なので Amdahl の法則における並列化の上限に効き、startup は worker 数に応じて伸びる固定費、compute は実際に並列で走る部分にあたる。
この四分割により、実行時間がどの区間に消えているかを worker 数ごとに追える。

startup の測り方には近似が含まれる。
\verb|Pool()| は子プロセスがグラフの unpickle を終える前に制御を返すため、プール生成の呼び出しが返った時刻は初期化の完了時刻ではない。
そこで、全 worker へ \verb|chunksize=1| で noop タスクをまき、それらが返るまでを起動バリアとした。
これが近似にとどまるのは、\verb|imap_unordered| が後続のバッチを、noop のまき方によっては初期化直後の特定の worker に偏って渡しうるためで、厳密な転送完了時刻とは一致しない。
この近似の性格は、startup の絶対値を手法間で比較するときの限界として第5章でも断る。

compute 区間では、worker ごとの稼働時間も同じ実行から取得する。
各バッチ処理が自プロセスの PID と所要時間を返し、親が PID ごとに積算する。
これにより、最も忙しい worker と最も暇な worker の稼働時間の差、および遊休率を測れる。
遊休率は、全 worker の稼働時間の合計を「worker 数 $\times$ compute 時間」で割った値を1から引いて求める。
バッチを一つも受け取らなかった worker も処理能力としては存在するので、分母は稼働した worker 数ではなく起動した worker 数を使う。
この遊休率が、負荷分布の偏りが worker の遊びとして現れる度合いを示す指標になる。

\todo{第5章（実験）で示す測定結果への具体的な相互参照（節番号・図表番号）を、実験章の確定後に補う}
